\documentclass[11pt]{article}
\usepackage[margin=1in]{geometry}
\usepackage[T1]{fontenc}
\usepackage[utf8]{inputenc}
\usepackage{amsmath,amssymb,amsthm,mathtools}
\usepackage{mathrsfs}
\usepackage{booktabs,tabularx,array,longtable}
\usepackage{enumitem}
\usepackage{hyperref}
\usepackage{microtype}
\usepackage{xcolor}
\usepackage{fancyhdr}
\usepackage{lastpage}
\hypersetup{colorlinks=true,linkcolor=blue,citecolor=blue,urlcolor=blue,pdftitle={From Search Dynamics to Semi-Ideal Computers},pdfauthor={Brian Tenneson}}

\newtheorem{definition}{Definition}[section]
\newtheorem{theorem}[definition]{Theorem}
\newtheorem{proposition}[definition]{Proposition}
\newtheorem{corollary}[definition]{Corollary}
\newtheorem{conjecture}[definition]{Conjecture}
\theoremstyle{remark}
\newtheorem{remark}[definition]{Remark}
\newtheorem{caution}[definition]{Caution}

\newcommand{\Qnn}{\mathbb{Q}_{\ge 0}}
\newcommand{\N}{\mathbb{N}}
\newcommand{\Fix}{\operatorname{Fix}}
\newcommand{\Cost}{\operatorname{Cost}}
\newcommand{\Ind}{\operatorname{Ind}}
\newcommand{\Dep}{\operatorname{Dep}}
\newcommand{\SIC}{\operatorname{SIC}}
\newcommand{\ATP}{\operatorname{ATP}}
\newcommand{\DATA}{\operatorname{DATA}}
\newcommand{\UNKNOWN}{\mathrm{UNKNOWN}}

\pagestyle{fancy}
\fancyhf{}
\lhead{Search Dynamics, SICs, and Rational Theorem Mazes}
\rhead{Tenneson, v0.1}
\setlength{\headheight}{14pt}
\cfoot{\thepage\ of \pageref{LastPage}}

\title{\textbf{From Search Dynamics to Semi-Ideal Computers}\\[0.4em]
\Large Cumulative Verified Histories, Rationally Weighted Theorem Mazes, and Reflective ATPs\\[0.6em]
\large A General Framework for ATPs, ALDs, and AMLDs, with DATA as a Distinguished Case}
\author{Brian Tenneson}
\date{August 2026 -- Version 0.1}

\begin{document}
\maketitle

\begin{center}
\textbf{Suggested arXiv classifications:} primary \texttt{cs.LO}; secondary \texttt{cs.AI}; cross-list candidate \texttt{math.CO}.
\end{center}

\begin{abstract}
This sequel develops a general bridge from staged verifier-backed search dynamics to semi-ideal computers (SICs), automated theorem provers (ATPs), automated logical deciders (ALDs), and automated metalogical deciders (AMLDs). The key construction is the \emph{cumulative history of verified information}: the permanent ledger obtained by unioning all mathematical information that has crossed an exact verification boundary through a given stage. The complete instantaneous search state may backtrack, restart, reallocate resources, forget heuristic proposals, and move nonmonotonically; its cumulative verified history is monotone by construction and therefore supplies the natural SIC state.

The paper then makes the graph-theoretic identification precise. Under a suitable encoding, theorem proving is weighted maze solving: vertices are proof or settlement states, directed edges are admissible search or inference transactions, exits are certificate-bearing states, and edge costs are exact nonnegative rational numbers. Any maze-search method whose correctness theorem holds for the relevant class of rationally weighted directed graphs can be pulled back to an ATP policy under the corresponding hypotheses. This yields a disciplined search-reduction program based on novelty DAGs, quotienting, conservative compilation, shared verified landmarks, incumbent bounds, branch-and-bound pruning, Dijkstra/A*-style search, and multi-agent navigation. The resulting theory does not claim that long proofs become short or that computational hardness disappears. Its objective is narrower: reduce the part of the proof ocean that must actually be searched before a known-to-exist certified fixed point becomes computationally accessible.

Finally, a sufficiently expressive induced ATP can be aimed at an effective code of the ATP, ALD, AMLD, or controller that generated it. This creates a verifier-gated meta-ATP able, within ordinary incompleteness and self-reference limits, to prove bounded, syntactic, relative, or conditional properties of the machine whose search is being optimized. DATA is treated as a distinguished implementation rather than as the general definition.
\end{abstract}

\tableofcontents
\clearpage

\section{Purpose and scope}
The motivating observation is that several strands of the earlier program fit into one diagram:
\[
\text{search dynamics}
\longrightarrow
\text{cumulative verified history}
\longrightarrow
\text{SIC}
\longrightarrow
\begin{cases}
\text{ATP},\\
\text{ALD},\\
\text{AMLD}.
\end{cases}
\]
A rich implementation such as $\DATA(x)$ is then a special case. The general mathematics should not depend on the names Professor, Counselor, Picard, Dreamer, Horizon, Prover, Refuter, or Independence. Those modules matter operationally, but the structural results concern any verifier-backed staged search process satisfying the stated hypotheses.

The central computational premise is equally modest. In many theorem-proving or conjecture-settling instances, a valid certificate and hence a terminal state exists mathematically, while available hardware may not reach it. The principal problem is therefore not fixed-point existence but \emph{computational accessibility}: how to shorten the hitting time, reduce the number of states that must be examined, compress repeated structure, and spend resources on branches that are more likely to reach a certificate.

\section{General verified transition systems}
\begin{definition}[Verified transition system]
A verified transition system is a tuple
\[
\mathfrak M=(X,\mathcal F,\pi,V,\mathcal R)
\]
where $X$ is a set of complete instantaneous states, $\mathcal F=\{F_j:X\to X\}_{j\in J}$ is a family of admissible state transformations, $\pi$ is a controller or scheduling policy selecting transformations, $V$ is an exact verification gate for mathematical deposits, and $\mathcal R$ is the space of typed verified records.
\end{definition}

A run is a sequence
\[
x_0,x_1,x_2,\ldots,\qquad x_{n+1}=F_{j_n}(x_n),
\]
with $j_n$ selected by $\pi$ or by an allowed nondeterministic/stochastic rule. The state $x_n$ may include queues, heuristic parameters, random seeds, learned representations, resource allocations, untrusted proposals, checkpoints, and verified mathematical memory.

\begin{caution}[The complete state need not be monotone]
Nothing requires $x_n\subseteq x_{n+1}$ or any analogous order relation. A practical searcher may discard a queue, restore a checkpoint, lower a beam width, change a representation, or restart. Monotonicity will be imposed only on the verified-history projection.
\end{caution}

\section{The cumulative history of verified information}
This construction is the bridge from general search dynamics to SICs.

\begin{definition}[Verified stage output]
Let $V_n$ denote the set of mathematical records newly accepted by the independent verifier during the transition into stage $n$. Set $V_1=\varnothing$ by convention, so that stage 1 is exactly the input stage. Depending on the machine, later $V_n$ may include formulas, proved lemmas, proof certificates, refutation certificates, independence or dependence certificates, model witnesses, provenance data, formally checked derived rules, or other typed mathematical objects.
\end{definition}

\begin{definition}[Cumulative history of verified information]
For input $\Gamma$ and stage $n\ge 1$, define
\[
C_{\mathfrak M}(\Gamma,n)
=
\Gamma\cup\bigcup_{k=1}^{n}V_k.
\]
In particular, $C_{\mathfrak M}(\Gamma,1)=\Gamma$.
This is the \emph{cumulative history of verified information} through stage $n$.
\end{definition}

The phrase is meant literally. $C_{\mathfrak M}(\Gamma,n)$ is not a snapshot of everything the machine currently believes, proposes, scores, or searches. It is the permanent ledger of what has crossed the verification boundary. Search states may be abandoned; verified mathematics is not silently revoked.

\begin{proposition}[Cumulative-history monotonicity]
For every run and every $n$,
\[
C_{\mathfrak M}(\Gamma,n)
\subseteq
C_{\mathfrak M}(\Gamma,n+1).
\]
\end{proposition}
\begin{proof}
By definition,
\[
C_{\mathfrak M}(\Gamma,n+1)
=
C_{\mathfrak M}(\Gamma,n)\cup V_{n+1}.
\]
\end{proof}

\begin{remark}[Projection, not necessarily quotient dynamics]
The cumulative history is a monotone observation of a possibly nonmonotone machine. Two complete states can have the same cumulative verified history while differing in queues, resources, or controller state and therefore having different futures. Thus the cumulative history need not be Markov-sufficient. It is a canonical SIC state even when it is not by itself a deterministic dynamical state.
\end{remark}

\section{The cumulative-history SIC theorem}
The SIC framework used in \emph{Depths of a Simulation} consists of a stage-state map $C$ and halting map $H$ with an initial state, monotone accumulation, persistent halting, and freezing after halting \cite{depths}.

\begin{definition}[Certificate halting]
Fix a target condition $Z\subseteq\mathcal R$ of accepted terminal records. Define
\[
H_{\mathfrak M}(\Gamma,n)=1
\quad\Longleftrightarrow\quad
C_{\mathfrak M}(\Gamma,n)\cap Z\neq\varnothing.
\]
After the first such stage, the operational machine may halt. For SIC analysis, extend the halted computation constantly thereafter.
\end{definition}

\begin{theorem}[Cumulative-History SIC Theorem]
Suppose verified records are persistent and the computation is extended constantly after its first accepted terminal certificate. Then
\[
M_{\mathfrak M}=(C_{\mathfrak M},H_{\mathfrak M})
\]
is a semi-ideal computer on the chosen verified-record universe.
\end{theorem}
\begin{proof}
The initial condition is fixed by construction. Monotonicity is the previous proposition. If $H_{\mathfrak M}(\Gamma,n)=1$, then the accepted terminal record remains in every later cumulative history, so halting persists. By the constant extension convention, the cumulative state is frozen after the operational halt.
\end{proof}

\section{When the induced SIC is an ATP, ALD, or AMLD}
\subsection{ATP specialization}
If the verified records are formulas and proof certificates in a frozen formal system, and the machine halts on target $\varphi$ exactly when a verifier-accepted proof of $\varphi$ is in the cumulative history, then the induced SIC is a target-search ATP once the usual soundness and eventual-success conditions are imposed. The internal search machinery can be arbitrarily elaborate; the SIC sees only the monotone verified ledger.

\subsection{ALD specialization}
For a target sentence $\varphi$ over theory $G$, use typed terminal certificate classes
\[
\begin{aligned}
P&:G\vdash\varphi, & R&:G\vdash\neg\varphi,\\
I&:\Ind_G(\varphi), & \bot&:\text{certified inconsistency of the base presentation}.
\end{aligned}
\]
An ALD is then a typed cumulative-history SIC whose halting condition is the appearance of an accepted supported settlement certificate.

\subsection{AMLD specialization}
At the metalogical level, define
\[
\Dep_G(\varphi):=(G\vdash\varphi)\lor(G\vdash\neg\varphi),
\qquad
\Ind_G(\varphi):=\neg\Dep_G(\varphi),
\]
and type every meta-certificate by an explicit metatheory. An AMLD is the same SIC-lift principle applied one logical level higher. Timeout, stagnation, and budget exhaustion remain $\UNKNOWN$ unless a certificate language explicitly proves otherwise.

\section{Halting and absorbing fixed points}
Let $T:X\to X$ be a deterministic macro-update for a complete-state model and let $\mathcal S\subseteq X$ be the certified settlement states. Extend terminal states as absorbing:
\[
x\in\mathcal S\Longrightarrow T(x)=x.
\]
If every nonterminal complete state changes under $T$, then
\[
\Fix(T)=\mathcal S.
\]
This is the fixed-point settlement pattern proved for DATA 3.0 under an explicit round-counter hypothesis \cite{data3}.

\begin{proposition}[Fixed-Point/Halting Correspondence]
Assume $\Fix(T)=\mathcal S$, and define $H(\Gamma,n)=1$ exactly when $x_n\in\mathcal S$. Then
\[
H(\Gamma,n)=1
\quad\Longleftrightarrow\quad
x_n\in\Fix(T).
\]
\end{proposition}

\begin{caution}[Existence is not reachability]
The identity $\Fix(T)=\mathcal S$ characterizes terminal states; it does not imply that a given orbit reaches one. A fixed point can exist while its hitting time exceeds every practical hardware, memory, or wall-clock budget. The search problem is therefore to reduce the cost of reaching a fixed point, not merely to prove that fixed points exist.
\end{caution}

\section{Theorem proving as rationally weighted maze solving}
The graph-theoretic viewpoint from \emph{Depths of a Simulation} can be strengthened into a transfer principle. Proof search is treated as directed graph traversal: proof states are vertices, legal inference transactions are edges, and target states are exits \cite{depths}.

\begin{definition}[Rationally weighted theorem maze]
A rationally weighted theorem maze is
\[
\mathcal G=(V,E,s,Z,w)
\]
where $(V,E)$ is a directed graph, $s\in V$ is the initial proof/search state, $Z\subseteq V$ is the set of certificate-bearing exits, and
\[
w:E\to\Qnn
\]
assigns each directed edge an exact nonnegative rational cost. The cost of a finite path $P=(e_1,\ldots,e_m)$ is
\[
w(P)=\sum_{i=1}^{m}w(e_i)\in\Qnn.
\]
\end{definition}

Rational weights are a deliberate restriction. They can be stored, added, and compared exactly. This avoids making a logical optimality theorem depend on floating-point ordering accidents. Integer transaction counts are the special case $w(e)\in\N$.

\begin{theorem}[Rational-Weighted Maze Transfer Principle]
Let $\mathcal A$ be a maze-search algorithm whose correctness theorem holds for a class $\mathscr C$ of effectively presented directed graphs with nonnegative rational edge weights. Suppose a theorem-proving or settlement problem can be encoded as a member of $\mathscr C$ so that:
\begin{enumerate}[label=(\roman*)]
\item vertices encode admissible proof/search states;
\item directed edges encode admissible transactions or inference steps;
\item target vertices are exactly states containing a verifier-accepted certificate of the requested type; and
\item path weight is the declared ATP search cost.
\end{enumerate}
Then applying $\mathcal A$ to the encoded theorem maze yields an ATP/ALD/AMLD search policy with the same route-finding and path-cost guarantee, subject to the same graph-theoretic hypotheses.
\end{theorem}
\begin{proof}
The encoding is a cost-preserving correspondence between maze paths and admissible search paths. Therefore any path returned by $\mathcal A$ pulls back to an admissible ATP path, and any optimality guarantee over the maze path weights transfers to the declared ATP search cost. Logical acceptance still requires the independent certificate verifier.
\end{proof}

\begin{remark}[Examples]
For finite nonnegative rationally weighted theorem mazes, Dijkstra's algorithm gives a minimum-cost route. A* gives the same conclusion under the usual admissibility conditions, with the standard implementation qualifications concerning consistency and node reopening. On a finite DAG, dynamic programming after a topological ordering gives minimum-cost routes. These are not new graph theorems; the contribution is the exact pullback to verifier-backed proof or settlement search.
\end{remark}

\begin{caution}[Search cost versus native proof cost]
A weighted theorem maze may price search transactions, macro rules, verifier calls, or other operations. That metric need not equal native proof length. If a banked lemma is transparently expanded before native proof cost is measured, the bank may dramatically reduce search cost while leaving the native minimum proof Horizon unchanged \cite{data3}. Every reported shortest-path claim must therefore name the metric and presentation.
\end{caution}

\section{Search-space reduction before learning}
The graph interpretation provides several exact or theorem-backed reductions before any learned heuristic is trusted.

\subsection{Cumulative verified-history collapse}
Many transient trajectories can lead to the same permanent verified ledger. If the objective is mathematical settlement rather than replaying every heuristic choice, states that differ only in irrelevant transient history can often share downstream mathematical information. Care is required because the cumulative history alone need not determine future controller behavior, but it is already a major conceptual reduction in what must be remembered as mathematical progress.

\subsection{Novelty DAG}
In the novelty state theorem graph, retain only nontrivial transitions that add new verified information. Strict inclusion then rules out directed cycles, producing a DAG \cite{depths}. Identity moves, rediscovery loops, and reversible syntactic wandering can be suppressed or represented outside the verified-state graph.

\subsection{Quotienting by verified equivalence}
When the formal system supplies a sound equivalence relation and the search objective is invariant under it, equivalent presentations can be quotiented or canonically oriented. This is a standard graph reduction, but in ATP design it can eliminate large families of syntactic duplicates. The quotient must be verifier-respecting; heuristic similarity is not enough.

\subsection{Conservative compilation and macro edges}
\emph{Depths of a Simulation} proves a conservative compilation result for finite provable workloads: verified derivations can be packaged as sound derived rules, lowering target transaction counts without changing theoremhood \cite{depths}. In graph language, a previously expensive verified path is replaced by a reusable macro edge. This is the mathematical justification for viewing a theorem library as a compression layer over repeated proof structure.

\subsection{Shared verified landmarks}
A shared bank can turn repeated rediscovery into one verified discovery followed by cheaper retrieval. DATA 3.0 gives a conditional repeated-substructure savings theorem and separately shows that transparent lemma reuse need not change the native proof Horizon \cite{data3}. For search, the bank changes the geometry; for native minimum proof cost, it need not.

\subsection{Incumbent bounds}
The Proof-Horizon program supplies an especially useful reduction. Once any valid proof or settlement certificate of cost $B$ is verified, the unknown optimum $c^*$ satisfies
\[
c^*\le B.
\]
For minimum-cost certification, all candidates with cost $\ge B$ are irrelevant to the question whether a cheaper certificate exists. Each improved incumbent shrinks the remaining bounded region \cite{horizon,data201}.

\begin{proposition}[Rational branch-and-bound pruning]
Let $\mathcal G$ be a theorem maze with nonnegative rational edge weights. Let $B$ be the cost of a verified incumbent route to an exit. For a reached vertex $v$, let $g(v)$ be the exact cost already incurred and let $h(v)\in\Qnn$ be a valid lower bound on the cost of every continuation from $v$ to an exit. If
\[
g(v)+h(v)\ge B,
\]
then no route through $v$ can improve the incumbent, so $v$ may be pruned from a search whose sole objective is to find a route of cost $<B$.
\end{proposition}
\begin{proof}
Every continuation through $v$ has total cost at least $g(v)+h(v)\ge B$.
\end{proof}

This proposition combines the Horizon incumbent idea with ordinary weighted-maze branch-and-bound. Even the trivial lower bound $h=0$ removes every partial path whose cost has already reached $B$; stronger admissible lower bounds remove more.

\section{Heuristic branches and where to look}
The Branch-Covering Theorem in \emph{Depths} says, roughly, that a proof-covering nondeterministic policy contains a branch able to shadow a shortest proof; heuristics then face the navigation problem of allocating enough computation to such branches \cite{depths}. This does not identify the good branch, but it transforms the intelligence problem into ranking, covering, diversification, and resource allocation.

For a theorem maze, a learned critic can estimate quantities such as
\[
\widehat d(v,Z),\qquad
\widehat P(\text{settlement within budget}\mid v),\qquad
\widehat Q(v,a),
\]
without ever acquiring proof authority. The verifier decides acceptance. The learned model decides where to spend the next unit of search.

A practical architecture can combine exact and heuristic components:
\[
\begin{aligned}
&\text{exact rational costs}+\text{incumbent bound}+\text{admissible lower bounds}\\
&\qquad +\text{learned ranking}+\text{portfolio diversity}+\text{verified bank}.
\end{aligned}
\]
The exact pieces prune safely; the learned pieces reorder or allocate the remaining work.

\section{Prior synthetic benchmark evidence, with disclaimers}
The technical companion to the Proof Horizon reported an $L^*=100$ synthetic dry tournament on 20 independently seeded finite ``oceans.'' Each ocean had an externally established shortest successful route of exactly 100 edges. General searchers were not told $L^*=100$ and were not given the hidden shortest route. The stopping rule was first verifier-accepted proof \cite{horizoncomp}.

\begin{table}[ht]
\centering
\small
\setlength{\tabcolsep}{4pt}
\begin{tabularx}{\textwidth}{>{\raggedright\arraybackslash}Xrrrrr}
\toprule
System & Finished & Optimum first & Median $L$ & Median $E$ & Median time (ms)\\
\midrule
Depths-F specialized control & 20/20 & 20/20 & 100 & 100 & 0.002\\
BFS & 20/20 & 20/20 & 100 & 10928.5 & 3.013\\
DFS & 20/20 & 4/20 & 115 & 1747.5 & 0.487\\
Beam-64 & 19/20 & 16/20 & 100 & 5546.5 & 2.664\\
Data-ATP Hilbert-only ablation & 20/20 & 6/20 & 115 & 1317.5 & 7.115\\
Data-ATP full dry stand-in & 20/20 & 11/20 & 100 & 375 & 30.747\\
DATA 2.0.1 fast wing & 20/20 & 7/20 & 115 & 121.5 & 28.570\\
DATA 2.0.1 Horizon certifier & 20/20 & 20/20 certified & 100 & 10928.5 & 3.037\\
\bottomrule
\end{tabularx}
\caption{Previously reported synthetic dry results at $L^*=100$; reproduced here only as architectural evidence and benchmark context, not as a production-performance claim.}
\end{table}

\paragraph{Required disclaimer.}
These numbers are \emph{not} production benchmarks of the full Data-ATP or DATA repositories, do not establish superiority on natural mathematics, and should not be read as a universal speedup theorem. They are executable stand-ins designed to separate intrinsic path length from navigation overhead and to test measurement, failure reporting, and architecture. In that dry setting, the fuller Data-ATP stand-in used far fewer median expansions than BFS while the separate Horizon certifier recovered/certified the known optimum at exhaustive-search cost. The result motivates, but does not validate, the larger search-reduction program \cite{horizoncomp}.

\section{A proposed rational-weight theorem-maze benchmark}
The next benchmark should test the graph-theoretic claim directly and should use exact rational weights rather than floating-point surrogates.

\begin{definition}[Rational Ocean benchmark family]
A benchmark instance is a finite or effectively bounded directed theorem maze with:
\begin{enumerate}
\item a planted or independently certified minimum settlement cost $L_w^*\in\Qnn$;
\item exact edge weights $a/b$ stored in reduced rational form;
\item a frozen verifier and certificate language;
\item adjustable width, dead-end density, alternative successful routes, repeated substructures, and misleading locality;
\item optional typed exits for $P$, $R$, $I$, and $\bot$;
\item a hidden shortest route unavailable to general searchers.
\end{enumerate}
\end{definition}

The benchmark should compare at least a rational Dijkstra baseline, rational A* with frozen admissible heuristics, BFS on unit-weight controls, beam/best-first variants, portfolio search, shared-bank ablations, and the relevant ATP/ALD/AMLD architecture. Every addition must preserve exact certificate checking.

Primary measurements should include verified settlement rate and type, first-settlement cost, certified minimum settlement cost where attempted, node/edge expansions, wall-clock time, peak memory, verifier time, bank deposits and retrievals, restarts, and the gap
\[
\Delta=\Cost(\text{first verified settlement})-L_w^*.
\]
For minimum-cost experiments, report separately the cost of discovery and the cost of certification.

\paragraph{Benchmark disclaimer.}
A rational-weight synthetic maze is a controlled model of proof-search geometry, not a claim that real theorem libraries have the same distribution. Success is meaningful only if ranking survives held-out changes in width, decoys, weight denominators, representation, and eventually frozen natural formal-mathematics benchmarks.

\section{Reflective ATPs: proving things about the prover}
Once the cumulative-history SIC satisfies ATP conditions, the induced ATP can itself be given a formal target language containing an effective code of a frozen machine $M$.

Let $M^{\#}$ be a coded or tagged copy of the machine, in the spirit of the tagging constructions in the author's \emph{Depths} line \cite{depths}. The meta-ATP can then search for derivations of statements such as
\[
\text{``every accepted $P$-certificate of $M^{\#}$ verifies,''}
\]
\[
\text{``the terminal states of $M^{\#}$ are absorbing,''}
\]
\[
\text{``this scheduler is fair on this finite benchmark,''}
\]
or bounded search facts such as
\[
\text{``no certificate of declared cost $<B$ occurs in this effectively finite layer.''}
\]
If such a theorem is independently verified, it may be deposited in a typed meta-level bank and used by a controller as mathematical information about the search process.

\begin{theorem}[Meta-ATP construction, schematic form]
Suppose the cumulative-history lift of $M$ is an effective ATP $A_M$, and the background formal theory effectively represents a frozen code $M^{\#}$ together with the relevant syntactic verification predicates. Then $A_M$ can search for every derivable encoded property of $M^{\#}$ expressible in that theory. Any accepted result is a theorem of the background formal system, not an oracle statement about semantic truth.
\end{theorem}

\begin{caution}[No unrestricted self-certification]
This construction does not evade the classical incompleteness, undecidability, or undefinability barriers associated with Goedel, Church, and Tarski \cite{godel,church,tarski}. It is most credible for bounded, syntactic, relative, conditional, and verifier-checkable properties. A machine proving a coded local fact about its search is not the same thing as a machine proving its unrestricted global soundness.
\end{caution}

\section{DATA as a distinguished case}
DATA 3.0 is a rich instance of the general framework. Its complete state includes object theory, target, metatheory, typed verified bank, multiple search frontiers, untrusted proposals, history, learned parameters, resource allocation, candidate/accepted certificates, settlement status, and a monotone round counter \cite{data3}. Its active subsystems are modeled as transformations whose compositions generate a transformation semigroup.

The general theory above interprets DATA in four simultaneous ways:
\begin{enumerate}
\item the instantaneous DATA process is a controlled transition system or IFS-like semigroup action;
\item its cumulative verified history is a SIC;
\item the $P$ route specializes to an ATP, while the combined $P/R/I$ routes form an ALD and the nested dependence/independence route forms an AMLD;
\item the resulting theorem/settlement graph is a rationally weighted maze once exact rational transaction costs are declared.
\end{enumerate}

The DATA-specific intelligence layer can therefore be understood as a collection of maze-navigation policies. Professor scores frontiers; Counselor challenges representations; Picard supplies feedback; Creativity and Dreamer generate candidate moves; Horizon certifies bounded regions; the verified bank adds reusable macro structure. None of these changes the verifier's semantics.

\section{A search funnel for computational accessibility}
The synthesis of the Horizon and Depths programs suggests the following funnel:
\[
\begin{array}{c}
\text{raw staged search dynamics}\\[2pt]
\downarrow\\[2pt]
\text{cumulative verified-history SIC}\\[2pt]
\downarrow\\[2pt]
\text{novelty / quotient / DAG reduction}\\[2pt]
\downarrow\\[2pt]
\text{conservative compilation and shared verified landmarks}\\[2pt]
\downarrow\\[2pt]
\text{verified incumbent bound }B\\[2pt]
\downarrow\\[2pt]
\text{rational branch-and-bound / Dijkstra / A* / exact graph methods}\\[2pt]
\downarrow\\[2pt]
\text{learned ranking, portfolios, reconnaissance, and resurvey}\\[2pt]
\downarrow\\[2pt]
\text{certificate-bearing fixed point.}
\end{array}
\]
The strongest claim is not that the fixed point becomes easy. It is that each justified transformation can remove a class of states or routes that do not need to be searched for the declared objective.

\section{Research conjectures and implementation questions}
\begin{conjecture}[Verified-history compression]
On broad formal-mathematics families, searching a canonicalized graph of cumulative verified histories requires asymptotically or practically fewer expansions than searching raw transient histories under matched correctness constraints.
\end{conjecture}

\begin{conjecture}[Compiled-landmark acceleration]
A verified shared bank plus conservative macro compilation reduces rational weighted hitting time on families with repeated substructure, after retrieval and verification overhead are charged exactly.
\end{conjecture}

\begin{conjecture}[Admissible learned lower bounds]
For some nontrivial theorem families, a learned function can be calibrated into a valid rational lower bound $h(v)$ often enough to improve branch-and-bound or A* search while preserving exact minimum-cost guarantees.
\end{conjecture}

\begin{conjecture}[Reflective search improvement]
A meta-ATP can prove bounded or conditional facts about a frozen source ATP/ALD/AMLD that, when deposited into the verified meta-bank, measurably reduce later settlement cost without altering the verifier or certificate semantics.
\end{conjecture}

Key implementation questions include whether cumulative-history canonicalization is cheap enough to pay for itself; how to define rational edge costs that remain meaningful across machines; how much of proof-state equivalence can be checked exactly; whether macro compilation should affect search cost, native proof cost, or both; and which meta-theorems are strong enough to improve scheduling rather than merely restating architecture.

\section{Conclusion}
The central proposal is that theorem proving and conjecture settlement can be studied as exact rationally weighted maze solving over verifier-backed state graphs. The internal search process may be nonmonotone, stochastic, learned, or self-reflective, while its cumulative history of verified information forms a monotone SIC. Under additional target-search conditions that SIC is an ATP; with typed proof/refutation/independence certificates it becomes an ALD; one logical level higher it becomes an AMLD.

The graph-theoretic perspective matters because it imports a large body of route-finding ideas subject to explicit hypotheses. For finite nonnegative rational weights, Dijkstra-style minimum-cost search is exact. With admissible lower bounds, A*-style and branch-and-bound pruning can reduce the region below an incumbent. Novelty DAGs eliminate verified-state cycles, conservative compilation replaces repeated verified paths by macro edges, and shared lemma banks turn repeated proof structure into reusable landmarks. These reductions do not solve undecidability or guarantee practical access to every fixed point. They narrow the ocean in which the machine must search.

The further reflective possibility is that the induced ATP can itself reason about a coded source ATP, ALD, AMLD, or controller. If such meta-theorems can be verified and fed back into search, the machine does not merely learn statistically from its trajectories; it can sometimes use proved facts about its own frozen search architecture. DATA is a distinguished laboratory for this general theory, not its definition.

\begin{thebibliography}{21}

\bibitem{moore}
Edward F. Moore.
\emph{The Shortest Path Through a Maze}.
In \emph{Proceedings of the International Symposium on the Theory of Switching, Part II}, pp. 285--292. Harvard University Press, 1959.

\bibitem{dijkstra}
Edsger W. Dijkstra.
A note on two problems in connexion with graphs.
\emph{Numerische Mathematik}, 1:269--271, 1959.
\url{https://doi.org/10.1007/BF01386390}

\bibitem{astar}
Peter E. Hart, Nils J. Nilsson, and Bertram Raphael.
A formal basis for the heuristic determination of minimum cost paths.
\emph{IEEE Transactions on Systems Science and Cybernetics}, 4(2):100--107, 1968.
\url{https://doi.org/10.1109/TSSC.1968.300136}

\bibitem{clrs}
Thomas H. Cormen, Charles E. Leiserson, Ronald L. Rivest, and Clifford Stein.
\emph{Introduction to Algorithms}, 4th ed.
MIT Press, 2022.
\url{https://mitpress.mit.edu/9780262046305/introduction-to-algorithms/}

\bibitem{aima}
Stuart Russell and Peter Norvig.
\emph{Artificial Intelligence: A Modern Approach}, 4th ed.
Pearson, 2020.
\url{https://aima.cs.berkeley.edu/}

\bibitem{floyd}
Robert W. Floyd.
Algorithm 97: Shortest path.
\emph{Communications of the ACM}, 5(6):345, 1962.
\url{https://doi.org/10.1145/367766.368168}

\bibitem{warshall}
Stephen Warshall.
A theorem on Boolean matrices.
\emph{Journal of the ACM}, 9(1):11--12, 1962.
\url{https://doi.org/10.1145/321105.321107}

\bibitem{rationalcoding}
Brian Tenneson.
\emph{A Closed-Form Surjection from N onto Q}.
Version 2.2, 2026. Live GitHub copy:
\url{https://github.com/btenneson/pub/blob/main/math.LO_Logic/A%20Closed-Form%20Surjection%20from%20N%20onto%20Q%20by%20Brian%20Tenneson%202.2.pdf}

\bibitem{depths}
Brian Tenneson.
\emph{Depths of the theorem-search program: graph-theoretic proof search, SICs, and automated theorem proving}.
Public 2026 research manuscript in the \emph{Depths} line; live GitHub copy:
\url{https://github.com/btenneson/pub/blob/main/cs.LO_Logic_in_Computer_Science/Depths_of_Induction_ML_Guided_Proof_Search_and_Formalized_Self_Awareness_Version_47.pdf}

\bibitem{horizon}
Brian Tenneson.
\emph{The Universal Proof-Horizon Operator: From MIU and the nth Well-Formed Formula to ZFC, Predator, and Automated Logical Deciders}.
Version 1.0, August 2026. Live GitHub copy:
\url{https://github.com/btenneson/pub/blob/main/cs.LO_Logic_in_Computer_Science/The_Universal_Proof-Horizon_Operator_by_Brian_Tenneson.pdf}

\bibitem{data201}
Brian Tenneson.
\emph{DATA 2.0.1: Proof-Horizon Architecture / From a Target Formula to a Verified Minimum-Cost Proof}.
August 2026. Live GitHub copy:
\url{https://github.com/btenneson/pub/blob/main/cs.LO_Logic_in_Computer_Science/DATA_2_0_1_Proof_Horizon_Architecture_FINAL.pdf}

\bibitem{horizoncomp}
Brian Tenneson.
\emph{A Technical Companion to My Next Generation Automated Theorem Prover}.
2026; includes the Proof-Horizon/Ocean synthetic benchmark material used here. Live GitHub copy:
\url{https://github.com/btenneson/pub/blob/main/cs.LO_Logic_in_Computer_Science/A_Technical_Companion_to_My_Next_Generation_Automated_Theorem_Prover_v1.0.pdf}

\bibitem{hilbert}
Brian Tenneson.
\emph{Hilbert Space-Filling Curve Theorem Search}.
Version 0.5.3, 2026. Live GitHub copy:
\url{https://github.com/btenneson/pub/blob/main/cs.LO_Logic_in_Computer_Science/Hilbert_Space_Filling_Curve_Theorem_Search_v0.5.3.pdf}

\bibitem{data3}
Brian Tenneson.
\emph{DATA 3.0: AMLD/IFS Semigroup Architecture, Settlement Horizons, and Conjecture-Settling Research Program}.
August 2026. Live GitHub research library:
\url{https://github.com/btenneson/pub/tree/main/cs.LO_Logic_in_Computer_Science}

\bibitem{godel}
Kurt Goedel.
Ueber formal unentscheidbare Saetze der Principia Mathematica und verwandter Systeme I.
\emph{Monatshefte fuer Mathematik und Physik}, 38:173--198, 1931.
\url{https://doi.org/10.1007/BF01700692}

\bibitem{church}
Alonzo Church.
An unsolvable problem of elementary number theory.
\emph{American Journal of Mathematics}, 58(2):345--363, 1936.
\url{https://doi.org/10.2307/2371045}

\bibitem{tarski}
Alfred Tarski.
Der Wahrheitsbegriff in den formalisierten Sprachen.
\emph{Studia Philosophica}, 1:261--405, 1935.

\end{thebibliography}

\end{document}

